\subsection{Combined Segment and Page Translation}
Segment pre-translation and page-table translation are composed in a fixed order.
The effective-address unit first computes an address value from the selected addressing
mode. A segment register is then selected by the access class or by an explicit
segment-qualified effective-address form. The selected segment produces the linear
address that is either used directly or passed to the paging stage.
Architectural pointer values used by ordinary code are the address values before this
segment stage. If software needs the linear address produced by segment pre-translation,
it uses \texttt{SEGLEA}; the resulting value is not the normal representation of a
portable pointer under the ABI.

\[
\begin{aligned}
\mathrm{EA} &\rightarrow \text{segment pre-translation} \rightarrow \text{linear address}\\
\texttt{PTCR.PE}=1:\quad \text{linear address} &\rightarrow \text{page walk} \rightarrow \text{memory-system address}\\
\texttt{PTCR.PE}=0:\quad \text{linear address} &\rightarrow \text{memory-system address}
\end{aligned}
\]

Instruction fetches use the code-segment context, stack accesses use the stack-segment
context, and ordinary data accesses use the data-segment context unless the effective
address explicitly selects another segment. The segment stage does not replace paging
as the memory-protection mechanism. Disabled segments perform no window check and no
base addition. Translated-window segments check an offset and add the segment base.
Bounds-only segments check an already-linear address without adding the base. In both
enabled modes, the checked byte-addressed window has a page-granular base and span.

Let \(x\) be the address produced by effective-address calculation, and let the selected
segment register define:

\[
\mathit{base} = \mathit{base\_page} \times 4096,\qquad
\mathit{span} = m \times 2^e \times 4096,\qquad
\mathit{limit} = \mathit{base} + \mathit{span}.
\]

The base, span, and limit are computed in an unsigned domain wide enough to represent
the full 64-bit linear-address space plus one. If \(\mathit{limit} > 2^{64}\), the
segment register image is invalid and any access through it reports \texttt{PAGE\_FAULT}.

\begingroup\footnotesize
\begin{longtable}{@{}p{1.15in}p{1.0in}p{1.55in}p{1.80in}@{}}
\toprule
\textbf{Mode} & \textbf{Segment State} & \textbf{Check} & \textbf{Linear Address}\\
\midrule
\endhead
disabled & \texttt{m=0} & no segment-window check & \(\mathit{linear}=x\)\\
translated window & \texttt{m!=0, b=0} & \(0 \le x < \mathit{span}\) & \(\mathit{linear}=\mathit{base}+x\)\\
bounds-only window & \texttt{m!=0, b=1} & \(\mathit{base} \le x < \mathit{limit}\) & \(\mathit{linear}=x\)\\
\bottomrule
\end{longtable}
\manualtablecaption{Segment Pre-Translation Modes}
\endgroup

The translated-window mode is useful when an address is interpreted as an offset inside
the selected segment. The bounds-only mode is useful when the effective address is already
a linear address but should still be constrained to the segment window. In both enabled
segment modes, a failed segment-window check is reported through the page-fault exception
class.

\begingroup\footnotesize
\begin{longtable}{@{}p{1.15in}p{1.05in}p{3.25in}@{}}
\toprule
\textbf{Segment Result} & \textbf{Paging State} & \textbf{Final Translation Behavior}\\
\midrule
\endhead
\(\mathit{linear}=x\) from disabled segment & \texttt{PTCR.PE=0} & The effective address is used directly as the memory-system address. PTE presence, permission, cache, and addressing-type attributes are not applied.\\
\(\mathit{linear}=x\) from disabled segment & \texttt{PTCR.PE=1} & The effective address is interpreted as a linear address. Canonical-address checking, page walk, PTE permissions, PTE cache policy, and PTE addressing type are applied.\\
\(\mathit{linear}=\mathit{base}+x\) from translated segment & \texttt{PTCR.PE=0} & The segment performs the window check and base addition. The resulting linear address is used directly by the memory system.\\
\(\mathit{linear}=\mathit{base}+x\) from translated segment & \texttt{PTCR.PE=1} & The segment first checks and translates the offset. The resulting linear address is then canonical-checked and translated through the page tables.\\
\(\mathit{linear}=x\) from bounds-only segment & either state & The segment constrains the already-linear address to the selected window. If paging is enabled, the checked linear address is then translated by the page tables; otherwise it is used directly.\\
\bottomrule
\end{longtable}
\manualtablecaption{Segment and Paging Composition}
\endgroup

When paging is enabled, PTE permissions are evaluated after segment pre-translation.
For example, an access may pass the segment-window check and still fault because the
leaf PTE is not present, is not writable for a store, is not executable for an instruction
fetch, or is not accessible at the current privilege level. Conversely, paging cannot
rescue an address that fails the enabled segment-window check, because the page walk is
not attempted for that access.

Privilege checks that do not depend on address translation are performed before page-fault
reporting. Address failures from the active translation stages, including enabled
segment-window checks, canonical-address checks when paging is enabled, page-table walks,
PTE permission checks, and stack-memory translation, are reported through the
\texttt{PAGE\_FAULT} vector.
